\subsection{Thermal Properties of Rocks}

The thermal behavior of a geological medium is described first through the temperature field. Temperature is a scalar quantity, and in a two-dimensional time-dependent problem it may be written as
\begin{equation}
T = T(x,y,t),
\end{equation}
where \(T\) is temperature, \(x\) and \(y\) are spatial coordinates, and \(t\) is time. In rocks, temperature may vary because of the regional geothermal gradient, local heat sources, fluid movement, surface heating and cooling, or transient thermal disturbances. This temperature field is central to thermoelasticity because changes in temperature produce thermal expansion or contraction and may therefore modify the stress state of the rock \cite{robertson1988,turcotte2014}.

One of the most important thermal properties of a rock is thermal conductivity, commonly denoted by \(k\). Thermal conductivity measures the ability of a material to conduct heat. In SI units, it is expressed in watts per metre per kelvin, \(\mathrm{W\,m^{-1}\,K^{-1}}\). A rock with a high thermal conductivity transfers heat easily through its mineral lattice . A rock with a low thermal conductivity resists the transfer of heat and can maintain stronger temperature gradients . The thermal conductivity of geological materials is controlled by their mineral composition, porosity, pore fluids, cracks, temperature, pressure, and fabric. Quartz-rich rocks, carbonate rocks, mafic rocks and clay-rich rocks may have different thermal conductivities due to different constituent minerals and internal structures \cite{robertson1988,mavko2009}.

Porosity and the filling of pores are of particular importance for the effective thermal conductivity of rocks. The conduction is usually dominated by the solid mineral framework, but pores and cracks may impede the heat transfer between grains.Air-filled pores in dry rocks tend to reduce the bulk thermal conductivity because air conducts heat poorly compared to most rock-forming minerals. If the pores are filled with water or brine, the effective thermal conductivity may be greater than the dry condition because water conducts heat better than air and improves the thermal contact within the pore system. However, the effect depends on pore geometry, connectivity, saturation, pressure, and scale. Therefore, the thermal conductivity measured on a dry laboratory sample may differ from the effective conductivity of the same rock under saturated in situ conditions \cite{robertson1988,wolff1982}.

Thermal conductivity can also be anisotropic. In layered, foliated, or highly fractured rocks, heat may be conducted more readily in one direction than another. For example, bedding planes, aligned minerals or preferred orientations of fractures may allow heat conduction parallel to the fabric to a greater extent than across it. In such cases, thermal conductivity is not fully represented by a single scalar value and may be treated more generally as a tensor property. This is important for thermoelastic wave propagation because directional variation in heat flow can influence the distribution of thermal strain and thermal stress within the medium \cite{robertson1988,turcotte2014}.

Another important property is the specific heat capacity, usually denoted by \(C_p\). Specific heat capacity measures the amount of heat required to raise the temperature of a unit mass of material by one degree. Its SI unit is \(\mathrm{J\,kg^{-1}\,K^{-1}}\). In rocks, \(C_p\) depends mainly on mineral composition and temperature. Materials with great heat capacity can absorb more heat with little temperature increase. Materials with low heat capacity are more easily changed temperature under the same heat input. In thermoelastic problems, heat capacity is important because it controls how rapidly a rock warms or cools when heat is supplied or removed \cite{robertson1988,turcotte2014}.

For heat transfer in a volume of rock, the volumetric heat capacity is often more useful than the specific heat capacity alone. It is written as
\begin{equation}
\rho C_p,
\end{equation}
where \(\rho\) is density and \(C_p\) is specific heat capacity. The product \(\rho C_p\) has units of \(\mathrm{J\,m^{-3}\,K^{-1}}\) and represents the amount of heat required to raise the temperature of a unit volume of rock by one degree. This property expresses the thermal inertia of the medium. A rock with a high volumetric heat capacity will resist a rapid change in temperature, while a rock with a low volumetric heat capacity will respond more rapidly to heating or cooling. In porous rocks, the effective volumetric heat capacity depends on the solid mineral framework and on the fluid occupying the pore space, if present \cite{robertson1988,wolff1982}.

Thermal diffusivity combines thermal conductivity and volumetric heat capacity into a single parameter that describes the rate at which temperature disturbances spread through a material. It is commonly written as
\begin{equation}
\label{eq:thermal-diff}
\alpha_T = \frac{k}{\rho C_p},
\end{equation}
where \(\alpha_T\) is thermal diffusivity, \(k\) is thermal conductivity, \(\rho\) is density, and \(C_p\) is specific heat capacity. The unit of thermal diffusivity is \(\mathrm{m^2\,s^{-1}}\). Physically, thermal diffusivity expresses the ratio between the ability of the material to conduct heat and its ability to store heat. A high value of \(\alpha_T\) means that a temperature disturbance spreads and smooths out relatively quickly. A low value means that heat diffusion is slower and temperature gradients may persist for a longer time \cite{robertson1988,turcotte2014}.

It is important to distinguish thermal diffusivity (\alpha_T) from the coefficient of thermal expansion which is often denoted by (\alpha). The two quantities have different physical meaning, although the symbols are the same.. Thermal diffusivity describes the spreading of temperature through the medium, whereas the coefficient of thermal expansion describes the tendency of the material to change size when temperature changes. For an isotropic solid, linear thermal strain has the form $\varepsilon_{\mathrm{th}} = \alpha\Delta T$ already introduced in Eq.~\eqref{eq:eps-thermal},
where \(\varepsilon_{\mathrm{th}}\) is thermal strain, \(\alpha\) is the linear coefficient of thermal expansion, and \(\Delta T\) is the temperature change. The coefficient \(\alpha\) has units of \(\mathrm{K^{-1}}\). This relation means that heating produces expansion and cooling produces contraction, provided that the material is free to deform \cite{boley1960,nowacki1986}.

In rocks, thermal expansion is controlled by mineral composition and internal structure. Different minerals expand by different amounts when heated. In a polymineralic rock such as granite, quartz, feldspar, and mica may have different thermal expansion behavior, so heating can produce internal mismatch between adjacent grains. If expansion is constrained, thermal strain may generate thermal stress. This is one of the central mechanisms of thermoelasticity. In porous or fractured rocks, part of the thermal deformation may be accommodated by pore deformation, crack opening, crack closure, or grain-boundary movement, so the effective response of the rock mass may differ from that of the intact mineral framework \cite{robertson1988,jaeger2007,nowacki1986}.

Thermal properties of rocks are not fixed constants under all conditions. Thermal conductivity, heat capacity, thermal diffusivity, and thermal expansion may vary with temperature, pressure, saturation, fracture state, weathering, and scale. Higher temperature could alter the properties of minerals and affect the formation of microcracks.
Increased confining pressure can close cracks and improve grain contacts affecting both thermal and elastic behavior. Saturation may affect heat storage and heat conduction because pore fluids contribute to both volumetric heat capacity and effective conductivity. At the rock-mass scale, bedding, foliation, fractures, and lithological layering can produce anisotropic and spatially variable thermal properties \cite{robertson1988,mavko2009,turcotte2014}.

For thermoelastic wave propagation, these thermal properties are important because they control the thermal part of the coupled process. Thermal conductivity determines how heat flows, heat capacity determines how much heat is stored, thermal diffusivity determines how quickly temperature disturbances spread, and thermal expansion determines how temperature change is converted into strain. These quantities provide the physical bridge between heat conduction and mechanical response. The next section therefore considers thermal expansion and thermal stress in more detail, because these processes explain how a temperature disturbance can create mechanical stresses in a constrained geological medium.